\documentclass[12pt]{article}
\usepackage[T1]{fontenc}
\usepackage[utf8]{inputenc}
\usepackage{lmodern}
\usepackage{textcomp}
\usepackage{enumitem}
\usepackage{geometry}
\geometry{margin=1in}
\begin{document}
\begin{center}
Answer Key - Physics - Graduate Level Assessment 16 \
\small (Answers are ordered exactly as in the question paper.)
\end{center}

\section*{Multiple Choice}
\begin{enumerate}[label=\arabic*.]
\item \textbf{Answer:} A
\\ \textit{Options:} A) 0.87 \ensuremath{\times} $10^6$ m/sec; B) 0.72 \ensuremath{\times} $10^6$ m/sec; C) 0.55 \ensuremath{\times} $10^6$ m/sec; D) 0.40 \ensuremath{\times} $10^6$ m/sec; E) 1.2 \ensuremath{\times} $10^6$ m/sec
\\ \textit{Note:} The correct answer is obtained by applying the photoelectric equation, which states that the kinetic energy of the emitted electrons equals the energy of the incident photons minus the work function of the material, leading to a calculated velocity of 0.87 x $10^6$ m/sec after determining the photon energy from the wavelength.
\item \textbf{Answer:} C
\\ \textit{Options:} A) Dye laser; B) Semiconductor laser; C) Gas laser; D) Chemical laser; E) Fiber laser
\\ \textit{Note:} Gas lasers operate by utilizing the electronic transitions of free atoms in a gas state, allowing them to emit coherent light based on the energy levels of these atoms.
\item \textbf{Answer:} A
\\ \textit{Options:} A) kinetic energy; B) number of electrons; C) potential energy; D) temperature of the nucleus; E) mass
\\ \textit{Note:} In both fission and fusion events, while mass and energy may convert and transform, the principle of conservation of energy ensures that the total kinetic energy of the system remains unchanged before and after the reaction, as the energy released or absorbed is accounted for in the kinetic energy of the resultant particles.
\item \textbf{Answer:} D
\\ \textit{Options:} A) It will quadruple.; B) It will increase by a factor of 2.; C) It will become zero.; D) It will be halved.; E) It will double.
\\ \textit{Note:} When the resistance R is doubled in a circuit powered by a constant voltage V, the power dissipated is given by the formula P = $V^2$ / R, so doubling R results in halving the power P.
\item \textbf{Answer:} B
\\ \textit{Options:} A) 15 ohms; B) 1.5 ohms; C) 2.25 ohms; D) 50 ohms; E) 0.15 ohms
\\ \textit{Note:} The resistance of the nichrome wire is calculated using the formula R = ρ(L/A), where ρ is the resistivity of nichrome (approximately 1.10 x 10^-6 ohm-meters), L is the length of the wire in meters (2.25 m), and A is the cross-sectional area in square meters (1.5 x 10^-6 m²), resulting in a resistance of 1.5 ohms.
\end{enumerate}

\section*{True / False}
\begin{enumerate}[label=\arabic*.]
\setcounter{enumi}{5}
\item \textbf{Answer:} True
\\ \textit{Options:} A) True; B) False
\\ \textit{Note:} This statement is true because kinetic energy is proportional to the square of the velocity (KE = 1/2 $mv^2$), so if the kinetic energy doubles, the velocity must increase by a factor of √2, and since momentum is directly proportional to velocity (p = mv), it increases by a factor of 2, resulting in momentum quadrupling when considering the relationship between kinetic energy and momentum.
\item \textbf{Answer:} False
\\ \textit{Options:} A) True; B) False
\\ \textit{Note:} The statement is false because neutrons, which are neutral particles found in the nucleus, do not influence the behavior of electrons, which are negatively charged and governed by electromagnetic forces in their discrete orbits around the nucleus.
\item \textbf{Answer:} True
\\ \textit{Options:} A) True; B) False
\\ \textit{Note:} Removing two protons and two neutrons from neon-20, which has an atomic number of 10 and a mass number of 20, results in a nucleus with an atomic number of 12 and a mass number of 18, corresponding to magnesium-18.
\item \textbf{Answer:} True
\\ \textit{Options:} A) True; B) False
\\ \textit{Note:} The statement is true because the seesaw achieves equilibrium when the torques produced by the weights of big brother and little sister, which depend on their respective masses and distances from the pivot point, are equal, allowing them to balance effectively.
\item \textbf{Answer:} False
\\ \textit{Options:} A) True; B) False
\\ \textit{Note:} The statement is false because the minimum frequency required for positron-electron pair production is actually 1.022 MeV in energy, which corresponds to a frequency of approximately 2.45 x $10^20$ Hz, thus exceeding the value given in the statement.
\end{enumerate}

\section*{Long Form Response}
\begin{enumerate}[label=\arabic*.]
\setcounter{enumi}{10}
\item \textbf{Answer:} In the context of additive color theory, color mixing occurs when different wavelengths of light are combined, producing new colors perceived by the human eye. Specifically, when red light, which has a longer wavelength, overlaps with blue light, which has a shorter wavelength, the resultant mixture creates magenta. This is because the additive process involves the combination of the three primary colors of light: red, green, and blue. In this scenario, the absence of green light allows the combination of red and blue wavelengths to stimulate the corresponding cone cells in the retina, leading to the perception of magenta. Thus, magenta can be understood as the visual result of red and blue light overlapping without the presence of green, demonstrating the principles of additive color mixing.
\\ \textit{Note:} correct because it accurately describes additive color theory, illustrating how the combination of red and blue light, without green, results in the perception of magenta due to the additive mixing of the primary colors of light.
\item \textbf{Answer:} To calculate the horsepower developed by a jet aircraft engine producing a thrust of 3000 lb at a velocity of 600 mi/hr, we can use the formula: horsepower (hp) = (thrust (lb) \ensuremath{\times} velocity (mi/hr)) / 550. Substituting the given values, we find hp = (3000 lb \ensuremath{\times} 600 mi/hr) / 550, which results in approximately 6000 hp. This calculation illustrates the significant power output required for jet propulsion, highlighting the engine's efficiency in converting fuel energy into thrust. The high horsepower indicates the necessity for advanced engineering to optimize performance and fuel consumption in jet engines, crucial for improving the overall efficiency of air travel.
\\ \textit{Note:} correct because it accurately applies the horsepower formula, demonstrating the relationship between thrust, velocity, and power output, while also emphasizing the engine's efficiency in converting fuel energy into thrust at a high operational power level.
\end{enumerate}

\vfill
\noindent\textit{End of Answer Key}
\end{document}